\section{Discussion and limitations}\label{sec:limits}
\paragraph{Practical reading.}
A benchmark score that was also used to choose among $K$ candidates is an upper estimate whose optimism is of order $\sigma\sqrt{(1-\rho)/n}\sqrt{2\ln K}$: several points for $n\le200$ even with 16--32 candidates. Recommendations, each tested above: (1) report $K$ and $n$ with the score; (2) when questions are scarce, hold out a share (20--50\%) of the benchmark \emph{after} fixing the search and report its score and interval, valid for any optimiser including a human in the loop; (3) use max-$T$/Bonferroni bounds if a guarantee over the candidate set is needed, accepting conservativeness; (4) do not compare systems with different tuning budgets on the tuning benchmark; (5) do not trust bootstrap-optimism or empirical-Bayes corrections unless no candidate is clearly better.

\paragraph{Limitations.}
\emph{Models and tasks.} Two small models (the 0.5B is near chance) and one benchmark (MMLU-Pro, short questions) with letter-logit scoring rather than generation or chain-of-thought; larger models have different $\rho$ and spread, so the quantitative gaps (4--5 pts) should not be transferred. The formulas need only $\sigma$, $\hat\rho$, $n$, $K$, measurable from cached outcomes.
\emph{Search space and optimisers.} Candidates are random draws from a small grid of five knobs (true sd between candidates 1.5--2.7 pts), searched exhaustively or by successive halving. DSPy/MIPRO-style optimisers, whose proposals are adaptive and potentially far more numerous, were not run; Prop.~\ref{prop:bound} and split reporting apply to them unchanged, the numbers in Table~\ref{tab:curse} do not, and their effective $K$ is unknown, itself an argument for split reporting.
\emph{Simulation design.} The pool (1000/1500 questions) is the population and benchmarks are drawn with replacement; \S\ref{sec:fresh} and App.~\ref{app:oos} check this against disjoint questions and the without-replacement analogue, which bracket the pool estimate; the decision experiments use pool truth only. Replicate counts give s.e.\ of 0.05--0.1 pt for gaps; coverages are accurate to about $\pm1.5$ pts. Tests use the normal approximation for paired differences.
\emph{Protocols.} Max-$T$ and the Gaussian prediction assume a fixed candidate set and large $n$; the Gaussian prediction uses the population covariance and validates the mechanism, not a practical estimator. The shift experiment uses one split whose direction is specific to our models. We did not study repeated adaptive use of one hold-out, multiple benchmarks, or non-binary metrics. Pool size and $K_{\rm tot}$ differ between the models (1500/1000, 32/16) purely for CPU budget; no GPU or API was used.

\paragraph{Conclusion.}
Reporting the tuned score on the tuning questions is optimistic in a quantitatively predictable way, large relative to the differences practitioners care about, and able to reverse comparisons; a split report with a small selection share is a cheap, search-agnostic remedy with exact validity and a measurable price, whereas bootstrap-optimism and shrinkage corrections are not dependable.
